\vspace{-0.1in}
% In this paper, our model architecture is initialized from 
We constructed a model for visual instruction tuning based on MiniGPT4 \cite{zhu2023minigpt}, which consists of the Vision transformer \cite{liu2023covid} backbone as the image encoder, Vicuna \cite{chiang2023vicuna} as the text decoder and a pre-trained Q-Former to connect them. Vicuna is built upon LLaMA \cite{touvron2023llama} with stronger following ability. Following \cite{zhu2023minigpt}, the Q-Former is designed to extract visual features from the frozen image encoder. Before feeding into the frozen Vicuna as the visual prompt, we use a learnable linear projection layer to narrow the gap between extracted visual features with Vicuna embeddings. The detailed prompt is as follows: 

\textit{"Give the following image: <Img>ImageContent</Img>. You will be able to see the image once I provide it to you. Please answer my questions. \#\#\#Human: <Img><ImageHere></Img> Instruction \#\#\#Assistant: Answer"},

where \textit{"<Img>ImageContent</Img>"} is a template to let Vicuna know image features will be placed within the two special tokens \textit{"<Img>", "</Img>"}. \textit{"<ImageHere>"} is the actual slot for image features. \textit{"Instruction"} and \textit{"Answer"} represent the human instructions and model output text, respectively.